\documentclass[a4paper,10pt]{article}
\usepackage[utf8]{inputenc}
\usepackage{graphicx}
\usepackage{subfigure}
\usepackage{fancyhdr}
\usepackage{listings}
\usepackage{ gensymb }


%opening

\pagestyle{fancy} %eigener Seitenstil
\fancyhf{} %alle Kopf- und Fußzeilenfelder bereinigen
\fancyhead[L]{SEP Abgabegespräch SS2015/16} %Kopfzeile links
\fancyhead[C]{} %zentrierte Kopfzeile
\fancyhead[R]{\includegraphics[width=1cm]{../../template/tugraz_bw}} %Kopfzeile rechts
\renewcommand{\headrulewidth}{0.4pt} %obere Trennlinie
\fancyfoot[C]{} %Seitennummer
\renewcommand{\footrulewidth}{0.4pt} %untere Trennlinie


\lstset{language=C++, basicstyle=\ttfamily\tiny}  

\begin{document}

\section*{Extended Array}

Schreibe ein einfaches Programm um eine Klasse zu modellieren, die Bitoperationen auf ein Integerarray machen kann. 

Die Klasse soll ein Integerarray und dessen Größe im Konstruktor übernehmen. Es soll möglich sein den Bitoperator \&= oder \& auf die Klasse aufzurufen. Dieser Operator soll ein weiteres Objekt der Klasse oder einen Integer übernehmen. Im ersten Fall werden die Zahlen der Arrays elementweise behandelt (x = 1 $\rightarrow$ lengthOfArray : array1[x] \& array2[x]). Wenn die Längen der 2 Arrays nicht gleich sind, soll 'dimension mismatch' ausgegeben werden. Im zweiten Fall wird für jedes Element der Integer verwendet. 

\subsection*{Beispiel}
\begin{verbatim}
int main() {

    int arr[] = {1, 2, 3, 4, 5};
    int arr2[] = {0, 0, 0, 0, 0, 0};

    BitArray a(arr, 5);
    BitArray b(arr2, 6);

    a &= b;
    a &= 1;
    
    for(int i = 0; i < 5; i++) {
      cout << a.array_[i] << endl;
    }
    return 0;
}
\end{verbatim}
Output: 
\begin{verbatim}
dimension mismatch
1
0
1
0
1
\end{verbatim}

\subsection*{Aufruf}
\texttt{./bitarray}

\subsection*{Bewertung}
Dokumentation \& Programmierstil,  Dokumentation \& Programmierstil, und Struktur \& Korrektheit (je 1 Punkt)
\newpage


\end{document}
